\chapter{Power-Up and the Outside World}
\label{ch:power-up}

\section{Starting the System}
\label{sec:system-start}

When power is applied, the reset circuit holds \activelow{RSTI} asserted while the supplies and the clock settle. The
CPU, and through \activelow{RSTO} the two other chips, are held in reset. When \activelow{RSTI} is released the CPU
holds \activelow{RSTO} for a further six cycles, and then it fetches its reset vector and starts the program
(Section~\ref{sec:reset}). How long the reset circuit holds \activelow{RSTI}, and what it waits for, are not yet
specified.

The program then has a console in the following state, and it sets up each part before it uses it:

\begin{table}[H]
	\centering
	\begin{tabular}{lp{8.4cm}}
		\toprule
		\textbf{Part}  & \textbf{State at the start of the program}                                                                                                                                   \\
		\midrule
		CPU            & At the address in the reset vector, by convention the system ROM's start-up routine, with its registers as the architecture defines them (\archmanual{}, exceptions chapter) \\
		SIRCADA        & Blank, in Mode 1, every writable register zero except those the registers chapter lists (\videomanual{})                                                                     \\
		SIRCAS         & Silent: the output is disabled, the registers hold their reset values, and the audio memory holds nothing the program can rely on (\audiomanual{})                           \\
		Working memory & Not yet specified                                                                                                                                                            \\
		\bottomrule
	\end{tabular}
	\caption{State After Reset}
	\label{tab:reset-state}
\end{table}

At reset the CPU fetches its reset vector through the boot overlay, from the cartridge, or from the system ROM if there
is no cartridge (Section~\ref{sec:boot-overlay}). By convention the vector holds the address of the system ROM's
start-up routine. The routine clears the overlay, writes the default vectors into the working memory, sets the system
up and then jumps to the program on the cartridge. A game may set the vector to its own boot code instead. What the
routine sets up, and the address it jumps to, are not yet specified.

\subsection{Without a Cartridge}

If there is no cartridge, the system ROM answers at the low words, and its own vector table sends the CPU to a routine
that shows the SIRCUS logo and an icon to say that the cartridge is missing. If a controller button is held down when
the console is reset, the routine instead enters a simple monitor, which can load a program over the serial port. The
button, the format and speed of the load, and where the program is placed in the working memory are not yet specified.

\subsection{Order of Setting Up}

The three chips do not depend on each other, so a program may set them up in any order. Two rules apply to each:

\begin{itemize}
	\item \textbf{SIRCADA.} Install the interrupt handlers before setting the interrupt enable bits, because the vertical
	      sync interrupt cannot be masked (\videomanual{}, timing chapter).
	\item \textbf{SIRCAS.} Write the audio memory first, then the registers, and set the output enable last
	      (\audiomanual{}, programming chapter). Load the audio memory in vertical blanking if the picture is running, or
	      with the picture blanked, since the window is slow (Section~\ref{sec:reaching}).
\end{itemize}

\section{The Video Output}

SIRCADA produces an analogue colour signal on three outputs, R, G and B, and a composite sync, \activelow{CSYNC}. An
output reference, IREF, sets the full-scale colour level. During reset the three outputs are at black (\videomanual{},
package chapter). A video encoder turns them into composite video (Section~\ref{sec:display-wiring}), which leaves the
console on the video jack. The encoder and the reference are not yet specified.

\section{The Audio Output}

SIRCAS produces a serial stream of signed 16-bit stereo samples, and the TDA1543 turns it into two analogue currents
(Section~\ref{sec:converter}). The output stage (Section~\ref{sec:output-stage}) turns each current into a voltage and
filters it, and the board needs a filter that removes the images of the signal around the sample rate (\audiomanual{},
timing chapter), which that filter does in part. The signals leave on the left and right jacks
(Section~\ref{sec:av-connector}).

\section{The Audio and Video Connector}
\label{sec:av-connector}

The sound and the picture leave the console on three RCA jacks, in the usual colours:

\begin{table}[H]
	\centering
	\begin{tabular}{lll}
		\toprule
		\textbf{Jack} & \textbf{Signal} & \textbf{Level}                                                 \\
		\midrule
		Yellow        & Composite video & 1 V peak to peak into 75 $\Omega$                              \\
		White         & Left audio      & About 1.1 V RMS at full scale (Section~\ref{sec:output-stage}) \\
		Red           & Right audio     & About 1.1 V RMS at full scale                                  \\
		\bottomrule
	\end{tabular}
	\caption{The Audio and Video Jacks}
	\label{tab:av-jacks}
\end{table}

\noindent A television or a monitor with composite and audio inputs connects directly. An optional adapter, which is not part of
the console and is not described in this manual, takes the three signals and modulates them onto a television channel
for a set that has only an aerial input.

\section{Other Connections}

The serial port leaves the console on a connector. It links the console to another console, to a modem or to a
debugger. Its levels are described in Section~\ref{sec:serial-device}. The connector, the handshake lines and the speed
range are not yet specified. The controller ports are described in Section~\ref{sec:controller-ports}; their connectors
are not yet specified. The cartridge connector and the physical form of the console are not yet specified either.

\section{Power}
\label{sec:power}

Each chip has its own supply pins: the CPU has two pairs of power and ground pins, SIRCADA has two pairs and a separate
analogue supply, AVCC and AGND, and SIRCAS has power and ground pins whose number is not yet specified. The supply
voltages of those chips, the regulation, the decoupling and the power connector are not yet specified.

The converter runs from the same $+5$~V as the logic, and needs no other supply:

\begin{table}[H]
	\centering
	\begin{tabular}{llrrr}
		\toprule
		\textbf{Pin} & \textbf{Supply} & \textbf{Range} & \textbf{Typical current} & \textbf{Maximum current} \\
		\midrule
		5            & $+5$ V          & 3.0 to 8.0 V   & 50 mA                    & 60 mA                    \\
		\bottomrule
	\end{tabular}
	\caption{Converter Supply}
	\label{tab:converter-supplies}
\end{table}

\noindent The part dissipates about 250~mW. It rejects 100~Hz ripple on its supply by only 50~dB, so it takes a branch of
the $+5$~V supply of its own, decoupled with 100~nF at the pin. The output stage runs from the same $+5$~V. The serial device, the MC68681, takes $+5$~V
at up to 150~mA.
